\documentclass[12pt,a4paper]{article}

\usepackage[polish]{babel}
\usepackage[utf8]{inputenc}
\usepackage[T1]{fontenc}
\usepackage{geometry}
\usepackage{graphicx}
\usepackage{float}
\usepackage{amsmath}
\usepackage{listings}

\geometry{margin=2.5cm}

\lstset{
    basicstyle=\ttfamily\small,
    frame=single,
    breaklines=true,
    columns=fullflexible,
    keepspaces=true,
    showstringspaces=false
}

\begin{document}

\section*{Programowanie sieciowe zad 1B}

\textbf{Autorzy:} \\
Adrian Tarach (Lider) \\
Wojciech Kryński\\
Grzegorz Niespodziany \\
Data: 22.11.2026, Wersja: 1.0

\section*{Treść zadania}

Klient powinien wysyłać do serwera strumień danych w pętli (tzn. danych powinno być „dużo”, minimum rzędu kilku
MB). Serwer powinien odbierać dane, ale między odczytami realizować sztuczne opóźnienie (np. za pomocą funkcji
sleep()). W ten sposób symulujemy zjawisko odbiorcy, który „nie nadąża” za szybkim nadawcą. Stos TCP będzie
spowalniał nadawcę, aby uniknąć tracenia danych. Należy zidentyfikować objawy tego zjawiska po stronie klienta
(dodając pomiar i logowanie czasu) i krótko przedstawić swoje wnioski poparte uzyskanymi statystykami czasowymi.
Wskazane jest też przeprowadzenie eksperymentu z różnymi rozmiarami bufora nadawczego po stronie klienta (np. 100
B, 1 KB, 10 KB).


\section*{Opis rozwiązania}

\subsection*{Opis implementacji klienta (\texttt{sock\_client1.c})}

\textbf{1. Funkcja \texttt{write\_all()}} \\
Żeby mieć pewność, że cała porcja danych faktycznie zostanie wysłana, napisaliśmy funkcję pomocniczą \texttt{write\_all()}. Wywołuje ona \texttt{write()} w pętli \texttt{while (sent < len)}. Jeśli \texttt{write()} wyśle tylko część danych, dodajemy zwrócone \texttt{n} do licznika \texttt{sent}, przesuwamy wskaźnik w buforze (\texttt{buf + sent}) i w następnej iteracji wysyłamy resztę (\texttt{len - sent}). Gdy \texttt{write()} zostanie przerwane sygnałem (\texttt{n == -1} i \texttt{errno == EINTR}), ponawiamy próbę. Przy innym błędzie albo zerwanym połączeniu funkcja zwraca \texttt{-1}.

\begin{lstlisting}
static int write_all(int sock, const char *buf, size_t len) {
    size_t sent = 0;
    while (sent < len) {
        ssize_t n = write(sock, buf + sent, len - sent);
        if (n == -1) {
            if (errno == EINTR) {
                continue;
            }
            return -1;
        }
        if (n == 0) {
            return -1;
        }
        sent += (size_t)n;
    }
    return 0;
}
\end{lstlisting}

\textbf{2. Pomiar czasu zapisu} \\
Czas mierzymy bezpośrednio przed i po każdym wywołaniu \texttt{write\_all()} w głównej pętli \texttt{send\_data()}. Korzystamy z zegara \texttt{CLOCK\_MONOTONIC}, żeby zmiany czasu systemowego nie wpływały na wyniki. Różnicę dwóch znaczników (\texttt{t1}, \texttt{t2}) przeliczamy na mikrosekundy (\texttt{dt}).

Pojedynczych pomiarów nie zapisujemy w pamięci. W każdej iteracji aktualizujemy tylko minimum (\texttt{min\_us}), maksimum (\texttt{max\_us}) i sumę (\texttt{sum\_us}), a na końcu liczymy z nich średnią i przepustowość. Do porównania różnych rozmiarów bufora to wystarczy.

\begin{lstlisting}
struct timespec t1, t2;
clock_gettime(CLOCK_MONOTONIC, &t1);
if (write_all(sock, buf, chunk) == -1) {
    perror("writing on stream socket");
    return -1;
}
clock_gettime(CLOCK_MONOTONIC, &t2);
double dt = (t2.tv_sec - t1.tv_sec) * 1e6 + (t2.tv_nsec - t1.tv_nsec) * 1e-3;
\end{lstlisting}

\textbf{3. Podział danych na porcje} \\
Klient alokuje bufor o rozmiarze \texttt{buf\_size} wypełniony znakami \texttt{'A'} i w pętli \texttt{while (sent < total\_bytes)} wysyła kolejne porcje. Rozmiar porcji to mniejsza z wartości: \texttt{buf\_size} albo liczba bajtów, które jeszcze zostały (\texttt{total\_bytes - sent}). Dzięki temu zawsze wysyłamy dokładnie 4 MiB ($4\,194\,304$ B), nawet gdy rozmiar bufora nie dzieli tej wartości bez reszty (np. przy 100 B).

\textbf{4. Czekanie na odbiór wszystkich danych (\texttt{shutdown})} \\
Po wysłaniu wszystkich danych klient nie wywołuje od razu \texttt{close()}. Najpierw zamyka kierunek wysyłania przez \texttt{shutdown(sock, SHUT\_WR)}, a potem czeka w \texttt{read()}, aż serwer odbierze wszystko i sam zamknie połączenie:

\begin{lstlisting}
shutdown(sock, SHUT_WR);
char end_byte;
while (read(sock, &end_byte, 1) == -1 && errno == EINTR) {
}
\end{lstlisting}
Dzięki temu kontener klienta nie kończy się za wcześnie i serwer ma czas doczytać końcówkę danych z bufora.

\subsection*{Opis implementacji serwera (\texttt{sock\_server\_gi1.c})}

Po \texttt{accept()} serwer wchodzi do funkcji \texttt{handle\_client()}, gdzie w pętli \texttt{while (1)} czyta dane wywołaniami \texttt{read()} po maksymalnie 4096 B (\texttt{BSIZE}). Przerwanie sygnałem (\texttt{EINTR}) powoduje ponowienie odczytu, a \texttt{read()} zwracające \texttt{0} oznacza, że klient zamknął swoją stronę (\texttt{shutdown}), więc wychodzimy z pętli. Po każdym udanym odczycie serwer śpi 5 ms (\texttt{nanosleep()}), co symuluje wolne przetwarzanie danych.

\begin{lstlisting}
while (1) {
    ssize_t rval = read(msgsock, buf, BSIZE);
    if (rval == -1) {
        if (errno == EINTR) {
            continue;
        }
        perror("reading stream message");
        break;
    }
    if (rval == 0) {
        printf("Ending connection\n");
        break;
    }
    if (delay_us > 0) {
        nanosleep(&req, NULL);
    }
}
\end{lstlisting}

\section*{Napotkane problemy i sposoby ich rozwiązania}

\textbf{Zawieszanie się drugiego testu (bufor 1024 B)}

Przy uruchamianiu trzech testów po kolei pierwszy (bufor 100 B) przechodził bez problemu, ale drugi (1024 B) blokował się po kilku pierwszych wywołaniach \texttt{write()} i nigdy się nie kończył.

Okazało się, że klient po wpisaniu wszystkich danych do bufora jądra od razu kończył działanie, a Docker usuwał jego kontener, zanim wolny serwer zdążył odebrać resztę danych i koniec połączenia. Serwer zostawał więc zablokowany na \texttt{read()} z pierwszego połączenia i nie wracał do \texttt{accept()}. Kolejny klient łączył się (połączenie czekało w kolejce \texttt{listen}), zapełniał bufory i stawał na \texttt{write()}.

Rozwiązaliśmy to, dodając w kliencie \texttt{shutdown(sock, SHUT\_WR)} po wysłaniu danych i czekanie w \texttt{read()}, aż serwer zamknie gniazdo. Od tej pory kontener klienta żyje, dopóki serwer nie doczyta wszystkiego, i wszystkie trzy testy przechodzą.


\section*{Opis konfiguracji testowej}

\begin{itemize}
    \item Kontenery komunikują się w wydzielonej sieci mostkowanej Dockera \texttt{z35\_network} przez swoje wirtualne interfejsy sieciowe.
    \item Serwer (\texttt{z35\_server}) dostaje dynamiczny adres IP z podsieci, wiąże gniazdo na \texttt{0.0.0.0} i nasłuchuje na porcie \texttt{8888/TCP}.
    \item Klient (\texttt{z35\_client}) dostaje kolejny adres z tej samej podsieci i łączy się po nazwie \texttt{z35\_server} na port \texttt{8888/TCP} z losowego portu lokalnego.
    \item Przesyłamy 4 MiB danych, kolejno dla buforów klienta 100 B, 1024 B i 10240 B. Opóźnienie na serwerze jest stałe: 5 ms po każdym wywołaniu \texttt{read()}, które zwraca do 4096 B.
\end{itemize}

\section*{Opis testowania i wyniki pomiarów}

Testy uruchamiamy skryptem \texttt{run.sh}, który buduje obrazy, uruchamia serwer w tle w sieci \texttt{z35\_network} i po kolei wywołuje klienta dla buforów 100 B, 1024 B i 10240 B. Poprawność działania sprawdziliśmy na podstawie wyników czasowych klienta i logów serwera, w których widać otwarcie i zamknięcie wszystkich trzech połączeń.

W tabeli ,,Czas całk.'' to suma czasów wszystkich wywołań \texttt{write\_all()} (bez czekania na zamknięcie połączenia), a przepustowość to 4 MiB podzielone przez ten czas. Każdy wariant uruchomiliśmy raz.

\begin{table}[H]
\centering
\footnotesize
\setlength{\tabcolsep}{4.5pt}
\begin{tabular}{|c|c|c|c|c|c|c|}
\hline
\textbf{Bufor klienta} & \textbf{Wywołania} & \textbf{Czas całk.} & \textbf{Przepustowość} & \textbf{MIN} & \textbf{AVG} & \textbf{MAX} \\
\hline
100 B & 41\,944 & 3,958 s & 1,011 MiB/s & 1,06 $\mu$s & 94,36 $\mu$s & 626\,144,66 $\mu$s \\
\hline
1024 B & 4\,096 & 4,027 s & 0,993 MiB/s & 1,38 $\mu$s & 983,24 $\mu$s & 435\,991,56 $\mu$s \\
\hline
10240 B & 410 & 3,307 s & 1,209 MiB/s & 1,76 $\mu$s & 8\,066,24 $\mu$s & 726\,118,70 $\mu$s \\
\hline
\end{tabular}
\end{table}


\section*{Wnioski}

\begin{itemize}
    \item Mimo stukrotnego zwiększenia bufora klienta (ze 100 B do 10240 B) czas wysyłania prawie się nie zmienił (3,31--4,03 s), a przepustowość wynosiła ok. 1,0--1,2 MiB/s. O tempie decyduje wolny serwer, a nie klient. Gdy bufor odbiorczy serwera się zapełnia, serwer ogłasza w TCP coraz mniejsze okno odbiorcze (aż do zera) i nadawca musi czekać. To jest sterowanie przepływem w TCP i dzięki niemu żadne dane nie giną.
    \item Zmierzona przepustowość jest wyższa niż to, co serwer jest w stanie odebrać (4096 B co 5 ms, czyli ok. 0,78 MiB/s, więc odczyt 4 MiB zajmuje mu co najmniej 5,1 s). Wynika to z tego, że mierzymy tylko czas \texttt{write()}. Gdy klient kończy pisać, ostatnie kilkaset KiB -- 1 MiB danych leży jeszcze w buforach jądra (nadawczym po stronie klienta i odbiorczym po stronie serwera), a serwer doczytuje je dopiero później.
    \item Czasy \texttt{write()} są bardzo nierówne. Najkrótsze zapisy trwały 1,06--1,76 $\mu$s, bo w buforze nadawczym było miejsce i dane były tylko kopiowane do jądra. Najdłuższe trwały 436--726 ms. Wtedy bufor był pełny i \texttt{write()} czekało, aż zwolni się w nim sporo miejsca (jądro nie budzi procesu po każdych kilku KB, tylko gdy zwolni się większa część bufora). To jest objaw ,,nienadążającego'' odbiorcy widoczny po stronie klienta.
    \item Większy bufor zmniejszył liczbę wywołań \texttt{write()} ze $41\,944$ do $410$. Średni czas jednego wywołania wzrósł z 94,36 $\mu$s do 8\,066,24 $\mu$s, ale to wynika głównie z tego, że prawie ten sam całkowity czas dzielimy przez ok. 100 razy mniej wywołań. Różnica w przepustowości (1,01 vs 1,21 MiB/s) jest niewielka, a każdy test był uruchomiony tylko raz, więc nie traktujemy jej jako istotnej.
\end{itemize}

\end{document}